% docs/justificacion.tex
\section{Justificación}

\subsection*{Por qué surge esta propuesta}

La necesidad de DIAPCE no nace de una carencia tecnológica abstracta, sino de una realidad concreta y repetida en el laboratorio de materiales y en el semillero de investigación de la UCC, sede Meta: cada cohorte de estudiantes parte prácticamente desde cero. Los registros de ensayos anteriores no se encuentran en un formato consultable, las condiciones de cada experimento se documentan de forma dispersa y los resultados acumulados por el grupo no se traducen en un saber colectivo que oriente las decisiones de dosificación siguientes. Cuando el conocimiento generado en prácticas sucesivas no se sistematiza, la curva de aprendizaje se reinicia con cada nueva generación de estudiantes, y la ventaja de pertenecer a un semillero activo se diluye.

A esto se suma el contexto ambiental del Meta: las condiciones de temperatura y humedad del territorio no son las que suponen las tablas genéricas de dosificación derivadas de normas internacionales. Diseñar mezclas de concreto estructural en esta región exige contrastar esas referencias con evidencia experimental propia, producida y validada localmente. Sin una herramienta que centralice y haga consultable ese acervo, cada decisión de dosificación se apoya en criterios individuales difíciles de comparar y reproducir.

\subsection*{Qué busca el proyecto}

DIAPCE busca convertir el conocimiento experimental acumulado por el semillero en un recurso colectivo, trazable y reproducible. No se trata de reemplazar el juicio del docente ni de automatizar el diseño de mezclas; se trata de ofrecer a los estudiantes una plataforma donde los datos de sus propios ensayos —resistencias medidas, condiciones ambientales registradas, aditivos empleados— se conviertan en la base real de las recomendaciones de dosificación.

En términos concretos, el proyecto busca: (i) centralizar los registros de diseño y los resultados de laboratorio en un repositorio estructurado y persistente; (ii) proporcionar asistencia en la selección de proporciones y aditivos a partir de los datos experimentales del semillero, con proyecciones de resistencia a distintas edades; y (iii) establecer una práctica de registro sistemático que permita comparar resultados entre sesiones, cohortes y condiciones ambientales distintas. Todo ello dentro de un entorno de uso exclusivamente formativo, sin pretensión de reemplazar herramientas de uso profesional o industrial.

\subsection*{Beneficios esperados}

Para el estudiante, DIAPCE representa un punto de contacto directo entre la teoría del diseño de mezclas y la evidencia que él mismo produce en el laboratorio. La herramienta hace visible el impacto de las variables ambientales y de la selección de aditivos sobre la resistencia del concreto, reforzando competencias que van más allá de la memorización de fórmulas: análisis crítico de resultados, formulación y contraste de hipótesis, y cultura de registro.

Para el semillero como grupo, el beneficio principal es la trazabilidad: la posibilidad de recuperar, comparar y discutir resultados de ensayos anteriores como si formaran parte de una misma conversación científica continua. Esto fortalece la identidad investigativa del grupo y eleva la calidad de las discusiones académicas entre cohortes.

Para los docentes, DIAPCE ofrece un instrumento de evaluación objetiva: los registros quedan almacenados con sus parámetros de entrada, lo que permite contrastar el proceso de decisión de cada estudiante con los resultados obtenidos y con el comportamiento del conjunto de datos regional.

\subsection*{Componente comunitario y aporte a la comunidad académica}

El valor de DIAPCE trasciende el uso individual. Al operar sobre un repositorio compartido de datos experimentales generados por los propios estudiantes de la UCC en el contexto del Meta, la herramienta construye —práctica a práctica— un patrimonio de conocimiento regional. Cada ensayo registrado enriquece la base desde la que se producen las recomendaciones futuras, lo que significa que el esfuerzo de investigación de una cohorte beneficia directamente a las siguientes.

Este modelo de acumulación progresiva del conocimiento es especialmente relevante en una región donde las condiciones ambientales particulares hacen insuficientes las referencias externas y donde la producción de datos locales tiene un valor formativo e institucional difícil de sustituir. DIAPCE propone que el semillero de investigación de la UCC deje de ser un espacio donde el conocimiento se genera y se pierde por ciclos, y se convierta en un espacio donde ese conocimiento se acumula, se discute y se refina de manera continua.

En síntesis, la justificación de este proyecto no descansa únicamente en la utilidad inmediata de la herramienta para cada usuario individual, sino en su capacidad de articular una comunidad de práctica: estudiantes, docentes y semilleros que comparten, contrastan y construyen sobre una misma evidencia experimental, en su propio contexto regional y bajo sus propias condiciones ambientales.
